% \subsection{Question answering}

% We find that previous retrieval-augmented LLMs
% tend to rephrase the retrieved facts without exploiting the knowledge of LLM itself. However,
% PoG enables LLM to synergistically infer
% from both the retrieved evidence graphs and its
% own knowledge. We attribute this ability to three
% aspects: 
% (1) Language Understanding, as LLM can
% comprehend and extract the knowledge from provided knowledge paths
% and the query in natural language, 
% (2) Knowledge Reasoning, as LLM can perform entity disambiguation and produce the final answer based on the provided finial knowledge paths, and (3) Knowledge Enhancement, as LLM can leverage its implicit knowledge to expand, connect, and improve the information relevant to the query. This ability is especially valuable when the external knowledge input is incomplete.

% Based on the three keys, we proposed our question answering method by 2 steps.
% \myparagraph{Path summarizing}

% Because some paths have too much text information, in some cases, even if the text information provided is wrong, LLM will answer the wrong question or provide the correct path, but the answer will be wrong. In order to solve this illusion problem, we proposed a strategy of summarizing the path first and then answering the question.

% Through the LLM prompt, by thinking about the provided path and split question, the triples related to the split question are extracted from the path and a new path is formed.
% The specific prompt used to guide the LLM in this selection process is detailed in the Appendix \textbf{TODO}.


% \myparagraph{Question Answering}
% we prompt LLMs to evaluate the provide answer if can answer all the split question first, then if can answer the original question.
% if both yes, the answer will be extraction from the provided path. all the process in this section is only call LLM once.

% if the node exploration reach the last depth and still no answer generated, we prompt LLM to answer the question by both provided knowledge and its inherited knowledge.The specific prompt used to guide the LLM in this selection process is detailed in the Appendix \textbf{TODO}.






\vspace{-1pt}
\subsection{Question Answering}
% This enhanced functionality is attributed to three main aspects:
% \begin{enumerate}
%     \item \textbf{Language Understanding}: LLMs can comprehend and extract knowledge from the provided knowledge paths and queries expressed in natural language.
%     \item \textbf{Knowledge Reasoning}: LLMs are capable of performing entity disambiguation and deriving the final answer from the final knowledge paths provided.
%     \item \textbf{Knowledge Enhancement}: LLMs can utilize their implicit knowledge to expand, connect, and refine the information pertinent to the query, which is invaluable when the external knowledge inputs are incomplete.
% \end{enumerate}
Based on the pruned paths in Section \ref{method:3pathprune}, we introduce a two-step question-answering method.

% \subsubsection{Path Summarizing}
% \myparagraph{Path summarizing}
% To address the issue of paths containing excessive or incorrect text information, which can lead to incorrect responses, we propose a strategy that involves summarizing the path before answering the question. Using an LLM prompt, the model thinks through the provided path and split question, extracting triples related to the split question to form a new, concise path. The specifics of the prompt used in this process are detailed in the Appendix \ref{appendix:prompt}.

% % \subsubsection{Question Answering}
% \myparagraph{Question answering}
% In this stage, we prompt the LLM to evaluate whether the provided answers can address all the split questions and then determine whether they can adequately respond to the original question. If both conditions are satisfied, the answer is extracted from the path provided and generated by prompting LLMs. 
% If the node exploration reaches the last permissible depth without generating an answer, we prompt the LLM to derive the answer using both the provided knowledge and its inherent knowledge. Details regarding the prompt used in this process are also specified in the Appendix \ref{appendix:prompt}.
\myparagraph{Path Summarizing}
To address hallucinations caused by paths with excessive or incorrect text, we develop a summarization strategy by prompting LLM to review and extract relevant triples from provided paths, creating a concise and focused path. Details of the prompts used are in Appendix \ref{appendix:prompt}.

% To address the issue of paths containing excessive or incorrect text information, which can lead to incorrect responses, we propose a strategy that involves summarizing the path before answering the question. Using an LLM prompt, the model thinks through the provided path and split question, extracting triples related to the split question to form a new, concise path. The specifics of the prompt used in this process are detailed in the Appendix \ref{appendix:prompt}.

\myparagraph{Question answering}
Based on the current reasoning path derived from path pruning and summarizing, we prompt the LLM to first evaluate whether the paths are sufficient for answering the split question and then the main question. If the evaluation is positive, LLM is prompted to generate the answer using these paths, along with the question and question analysis results as inputs, as shown in Figures \ref{fig:overall}. The prompts for evaluation and generation are detailed in Appendix \ref{appendix:prompt}. If the evaluation is negative, the exploration process is repeated until completion. If node expand exploration reaches its depth limit without yielding a satisfactory answer, LLM will leverage both provided and inherent knowledge to formulate a response. Additional details on the prompts can be found in Appendix \ref{appendix:prompt}.
